\subsection{Orthogonality Relation for Characters}

We keep the notation of Subsections 5 and 6. Recall that \(Z(K[G])\) consists of the central functions.

We define a symmetric bilinear mapping from \(K[G] \times K[G]\) to K by the formula

\[
\langle f, f ^ {\prime} \rangle_ {\mathrm{G}} = | \mathrm{G} | ^ {- 1} \sum_ {g \in \mathrm{G}} f _ {g} f _ {g ^ {- 1}} ^ {\prime}\tag{28}
\]

for every \(f = \sum f_{g}g\) and \(f' = \sum f_{g}'g\) belonging to K{[}G{]}. We have \(\langle f, f' \rangle_{\mathrm{G}} = |\mathrm{G}|^{-2}\tau (ff')\) .

PROPOSITION 4 (Orthogonality relation for characters)

For \(\lambda\) and \(\mu\) in \(\hat{\mathrm{G}}\) , we have \(\langle \chi_{\lambda},\chi_{\mu}\rangle_{\mathrm{G}} = \delta_{\lambda \mu}\) .

This is the specific case of relations (20) and (23) when the endomorphisms u and v are taken to be the identity.

COROLLARY. --- Let \(\pi\) and \(\pi'\) be finite-dimensional linear representations of G. In the field K, we have

\[
\langle \chi_ {\pi}, \chi_ {\pi^ {\prime}} \rangle_ {\mathrm{G}} = (\dim_ {\mathrm{K}} \operatorname{Hom} _ {\mathrm{G}} (\pi , \pi^ {\prime})) \cdot 1.\tag{29}
\]

We first suppose that \(\pi\) and \(\pi'\) are simple representations. The vector space \(\mathrm{Hom}_{\mathrm{G}}(\pi, \pi')\) has dimension 1 or 0 according to whether or not \(\pi\) and \(\pi'\) are isomorphic (Schur's lemma, VIII, p.~47, Proposition 2). In this case, formula (29) follows from Proposition 4.

In the general case, the representation \(\pi\) (resp. \(\pi'\) ) is the direct sum of simple representations \(\pi_{1},\ldots,\pi_{m}\) (resp. \(\pi_{1}',\ldots,\pi_{n}'\) ). The space \(\operatorname{Hom}_{\mathrm{G}}(\pi,\pi')\) is isomorphic to the direct sum of the spaces \(\operatorname{Hom}_{\mathrm{G}}(\pi_{i},\pi_{j}')\) for \(1\leqslant i\leqslant m\) , \(1\leqslant j\leqslant n\) , and we have

\[
\chi_ {\pi} = \chi_ {\pi_ {1}} + \dots + \chi_ {\pi_ {m}}, \quad \chi_ {\pi^ {\prime}} = \chi_ {\pi_ {1} ^ {\prime}} + \dots + \chi_ {\pi_ {n} ^ {\prime}}.
\]

By linearity, the proof of formula \((29)\) is reduced to the case of simple representations.

